% claim.tex -- AI4Math claim of record (AI-generated, 2026-09-28; instantiated
% from harness/templates/claim/claim.tex per SOP 08b).
% Machine-readable theorem anchors are the "% LEAN:" lines; scripts/paper-lint.sh
% and the claim audit check them. All Lean listings are verbatim, byte-identical
% contiguous excerpts of frozen artifacts of tasks/20260927-mossad-pend:
%   statements : formalized/01-pend-statements.lean (113 lines, sha256 78bccc97...)
%   proof file : final/01-pend-proved.lean         (259 lines, sha256 be3f31cb...)
% Line ranges used: statement 27-68, 71-89, 92-106, 109-112; final 1-259 (whole
% file). The frozen statement file's five literal "sorry" occurrences (four proof
% placeholders plus one header-comment mention) all lie outside the excerpts, as
% required by the paper-lint listing gate. Verification: audit/listing-verify.txt
% (audit/inspect-tex.py --verify).
\documentclass[11pt]{article}

\usepackage[T1]{fontenc}
\usepackage[utf8]{inputenc}
\usepackage{amsmath,amssymb,amsthm}
\usepackage{graphicx}
\usepackage{hyperref}
\usepackage{xurl}
\setlength{\emergencystretch}{3em}
% lstlean.tex — Lean style for the listings package.
% Lineage: Jeremy Avigad (2015), adapted from Assia Mahboubi's SSR style;
% inspired by Olivier Verdier's unixode.sty for the Unicode literate table.
% No CTAN package or upstream repo exists; papers carry this file in their
% source package (cap set arXiv:1907.01449, sphere eversion arXiv:2210.07746).
% AI4Math adaptation (AI-generated, 2026-09-26): sorry/admit/sorryAx elevated
% to a dedicated error tier, matching the pipeline's 0-sorry constitution.
\usepackage{listings}
\usepackage{xcolor}

\definecolor{keywordcolor}{rgb}{0.7, 0.1, 0.1}   % red keywords
\definecolor{tacticcolor}{rgb}{0.0, 0.1, 0.6}    % blue tactics
\definecolor{commentcolor}{rgb}{0.4, 0.4, 0.4}   % grey comments
\definecolor{symbolcolor}{rgb}{0.0, 0.1, 0.6}    % blue symbols
\definecolor{sortcolor}{rgb}{0.1, 0.5, 0.1}      % green sorts
\definecolor{errorcolor}{rgb}{0.6, 0, 0}         % dark red: sorry/admit
\definecolor{stringcolor}{rgb}{0.5, 0.1, 0.5}    % purple strings

\lstdefinelanguage{lean}{
  % Anything between $ becomes LaTeX math mode
  mathescape=false,
  % Comments may or not include Latex commands
  texcl=false,
  % keywords, list taken from lean-syntax.txt
  morekeywords=[1]{import, prelude, protected, private, noncomputable,
    definition, meta, renaming, hiding, parameter, parameters, begin,
    constant, constants, lemma, variable, variables, theory, print,
    theorem, example, open, section, namespace, universe, universes, end,
    modifier, modifiers, using, export, exposing, axiom, inductive,
    coinductive, with, structure, class, instance, attribute, local, where,
    by, have, show, let, in, if, then, else, match, do, for, fun, assume,
    from, take, calc, include, omit, infix, infixl, infixr, notation,
    macro, syntax, elab, set_option, initialize, deriving, opaque, abbrev,
    mutual, partial, scoped, termination_by, decreasing_by},
  % Sorts
  morekeywords=[2]{Sort, Type, Prop},
  % tactics, list taken from lean-syntax.txt (tactic keywords)
  morekeywords=[3]{intro, intros, apply, exact, refine, rfl, simp, simp_all,
    simpa, simp_rw, rw, rewrite, erewrite, ring, ring_nf, omega, decide,
    norm_num, norm_cast, induction, cases, rcases, obtain, constructor,
    ext, funext, conv, congr, field_simp, linarith, nlinarith, positivity,
    tauto, trivial, aesop, use, by_cases, by_contra, push_neg, contrapose,
    symm, trans, assumption, contradiction, exfalso, specialize,
    generalize, revert, clear, rename_i, subst, unfold, delta, change,
    convert, split_ifs, interval_cases, fin_cases, zify, gcongr,
    nth_rewrite, suffices, generalizing, at, only, native_decide},
  % warnings - constitution tier: must never survive into a final proof
  morekeywords=[4]{sorry, admit, sorryAx},
  literate=
    {α}{{\ensuremath{\alpha}}}1
    {β}{{\ensuremath{\beta}}}1
    {γ}{{\ensuremath{\gamma}}}1
    {δ}{{\ensuremath{\delta}}}1
    {ε}{{\ensuremath{\varepsilon}}}1
    {ζ}{{\ensuremath{\zeta}}}1
    {η}{{\ensuremath{\eta}}}1
    {θ}{{\ensuremath{\theta}}}1
    {ι}{{\ensuremath{\iota}}}1
    {κ}{{\ensuremath{\kappa}}}1
    {λ}{{\ensuremath{\lambda}}}1
    {μ}{{\ensuremath{\mu}}}1
    {ν}{{\ensuremath{\nu}}}1
    {ξ}{{\ensuremath{\xi}}}1
    {π}{{\ensuremath{\pi}}}1
    {ρ}{{\ensuremath{\rho}}}1
    {σ}{{\ensuremath{\sigma}}}1
    {τ}{{\ensuremath{\tau}}}1
    {υ}{{\ensuremath{\upsilon}}}1
    {φ}{{\ensuremath{\varphi}}}1
    {χ}{{\ensuremath{\chi}}}1
    {ψ}{{\ensuremath{\psi}}}1
    {ω}{{\ensuremath{\omega}}}1
    {Γ}{{\ensuremath{\Gamma}}}1
    {Δ}{{\ensuremath{\Delta}}}1
    {Θ}{{\ensuremath{\Theta}}}1
    {Λ}{{\ensuremath{\Lambda}}}1
    {Ξ}{{\ensuremath{\Xi}}}1
    {Π}{{\ensuremath{\Pi}}}1
    {Σ}{{\ensuremath{\Sigma}}}1
    {Φ}{{\ensuremath{\Phi}}}1
    {Ψ}{{\ensuremath{\Psi}}}1
    {Ω}{{\ensuremath{\Omega}}}1
    {ℵ}{{\ensuremath{\aleph}}}1
    {∀}{{\ensuremath{\forall}}}1
    {∃!}{{\ensuremath{\exists}!}}1
    {∃}{{\ensuremath{\exists}}}1
    {∈}{{\ensuremath{\in}}}1
    {∉}{{\ensuremath{\notin}}}1
    {∘}{{\ensuremath{\circ}}}1
    {∩}{{\ensuremath{\cap}}}1
    {∪}{{\ensuremath{\cup}}}1
    {⊆}{{\ensuremath{\subseteq}}}1
    {⊂}{{\ensuremath{\subset}}}1
    {⊇}{{\ensuremath{\supseteq}}}1
    {⊃}{{\ensuremath{\supset}}}1
    {⊄}{{\ensuremath{\nsubseteq}}}1
    {∅}{{\ensuremath{\emptyset}}}1
    {∧}{{\ensuremath{\wedge}}}1
    {∨}{{\ensuremath{\vee}}}1
    {¬}{{\ensuremath{\neg}}}1
    {⊤}{{\ensuremath{\top}}}1
    {⊥}{{\ensuremath{\bot}}}1
    {↔}{{\ensuremath{\leftrightarrow}}}1
    {→}{{\ensuremath{\rightarrow}}}1
    {←}{{\ensuremath{\leftarrow}}}1
    {↑}{{\ensuremath{\uparrow}}}1
    {⇒}{{\ensuremath{\Rightarrow}}}1
    {⇐}{{\ensuremath{\Leftarrow}}}1
    {⟹}{{\ensuremath{\Longrightarrow}}}1
    {↦}{{\ensuremath{\mapsto}}}1
    {≤}{{\ensuremath{\leq}}}1
    {≥}{{\ensuremath{\geq}}}1
    {≠}{{\ensuremath{\neq}}}1
    {≈}{{\ensuremath{\approx}}}1
    {≡}{{\ensuremath{\equiv}}}1
    {≃}{{\ensuremath{\simeq}}}1
    {≅}{{\ensuremath{\cong}}}1
    {×}{{\ensuremath{\times}}}1
    {·}{{\ensuremath{\cdot}}}1
    {±}{{\ensuremath{\pm}}}1
    {⊕}{{\ensuremath{\oplus}}}1
    {⊗}{{\ensuremath{\otimes}}}1
    {⋆}{{\ensuremath{\star}}}1
    {▸}{{\ensuremath{\blacktriangleright}}}1
    {⟨}{{\ensuremath{\langle}}}1
    {⟩}{{\ensuremath{\rangle}}}1
    {⦃}{{\ensuremath{\{\!\mid}}}1
    {⦄}{{\ensuremath{\mid\!\}}}}1
    {⟦}{{\ensuremath{[\![}}}1
    {⟧}{{\ensuremath{]\!]}}}1
    {⌊}{{\ensuremath{\lfloor}}}1
    {⌋}{{\ensuremath{\rfloor}}}1
    {⌈}{{\ensuremath{\lceil}}}1
    {⌉}{{\ensuremath{\rceil}}}1
    {√}{{\ensuremath{\surd}}}1
    {∣}{{\ensuremath{\mid}}}1
    {∤}{{\ensuremath{\nmid}}}1
    {∎}{{\ensuremath{\blacksquare}}}1
    {⋯}{{\ensuremath{\cdots}}}1
    {…}{{\ensuremath{\ldots}}}1
    {∑}{{\ensuremath{\sum}}}1
    {∏}{{\ensuremath{\prod}}}1
    {⁻¹}{{\ensuremath{^{-1}}}}1
    {¹}{{\ensuremath{^1}}}1
    {²}{{\ensuremath{^2}}}1
    {³}{{\ensuremath{^3}}}1
    {ⁿ}{{\ensuremath{^n}}}1
    {ᵐ}{{\ensuremath{^m}}}1
    {₀}{{\ensuremath{_0}}}1
    {₁}{{\ensuremath{_1}}}1
    {₂}{{\ensuremath{_2}}}1
    {₃}{{\ensuremath{_3}}}1
    {₄}{{\ensuremath{_4}}}1
    {₅}{{\ensuremath{_5}}}1
    {ₙ}{{\ensuremath{_n}}}1
    {ₘ}{{\ensuremath{_m}}}1
    {ₖ}{{\ensuremath{_k}}}1
    {ᵢ}{{\ensuremath{_i}}}1
    {ⱼ}{{\ensuremath{_j}}}1
    {ℤ}{{\ensuremath{\mathbb{Z}}}}1
    {ℕ}{{\ensuremath{\mathbb{N}}}}1
    {ℚ}{{\ensuremath{\mathbb{Q}}}}1
    {ℝ}{{\ensuremath{\mathbb{R}}}}1
    {ℂ}{{\ensuremath{\mathbb{C}}}}1
    {𝔽}{{\ensuremath{\mathbb{F}}}}1
    {𝒫}{{\ensuremath{\mathcal{P}}}}1
    {∗}{{\ensuremath{\ast}}}1
    {•}{{\ensuremath{\bullet}}}1,
  % Comments
  morecomment=[l]{--},
  morecomment=[s]{/-}{-/},
  morestring=[b]",
  stringstyle=\color{stringcolor},
  keepspaces=true,
  keywordstyle=[1]\color{keywordcolor}\bfseries,
  keywordstyle=[2]\color{sortcolor}\bfseries,
  keywordstyle=[3]\color{tacticcolor}\bfseries,
  keywordstyle=[4]\color{errorcolor}\bfseries\underline,
  escapeinside={(*@}{@*)},
  columns=fullflexible,
  basicstyle=\ttfamily\small,
  commentstyle=\color{commentcolor}\itshape,
  breaklines=true,
  showstringspaces=false,
  frame=single,
  framesep=2mm,
  xleftmargin=4mm,
  numbers=left,
  numberstyle=\tiny\color{commentcolor},
  numbersep=6pt,
}

% Inline Lean code in prose: \lean{Int.ModEq.pow}
\newcommand{\lean}[1]{\lstinline[language=lean,basicstyle=\ttfamily\normalsize]|#1|}

\lstset{language=lean}

% --- local rendering patch (2026-09-28) ---
% tectonic/XeTeX does not apply the listings literate table to multi-byte
% characters, so the Unicode set used by the listings of THIS note is mapped to
% math glyphs as active characters via the standard \lccode/\lowercase idiom
% (ASCII-only source, no extra packages). Restricted to symbols outside xeCJK's
% own punctuation table: 00B7 and 2014 also occur in the frozen sources but are
% special punctuation that xeCJK sets up at load time -- making them active first
% aborts the run with "Control sequence ... already defined" (observed 2026-09-27);
% they route through xeCJK + SimSun instead, as do all CJK ideographs and
% fullwidth/CJK punctuation (3001, 3002, 300C, 300D, 3010, 3011, FFxx).
% 21A6 (mapsto) and 27FA (iff) are new relative to the chorded-cycle-mod4 note
% (they occur in the pend sources and are not xeCJK punctuation). Rendering only:
% listing source bytes are untouched, so the byte-exactness checks are unaffected,
% and the \ifdefined guard leaves the pdflatex path identical to the template.
\ifdefined\XeTeXversion
  {\lccode`\~="2115 \lowercase{\global\catcode`~=\active \global\def~{\ensuremath{\mathbb{N}}}}}
  {\lccode`\~="2124 \lowercase{\global\catcode`~=\active \global\def~{\ensuremath{\mathbb{Z}}}}}
  {\lccode`\~="2190 \lowercase{\global\catcode`~=\active \global\def~{\ensuremath{\leftarrow}}}}
  {\lccode`\~="2192 \lowercase{\global\catcode`~=\active \global\def~{\ensuremath{\rightarrow}}}}
  {\lccode`\~="2194 \lowercase{\global\catcode`~=\active \global\def~{\ensuremath{\leftrightarrow}}}}
  {\lccode`\~="21A6 \lowercase{\global\catcode`~=\active \global\def~{\ensuremath{\mapsto}}}}
  {\lccode`\~="2200 \lowercase{\global\catcode`~=\active \global\def~{\ensuremath{\forall}}}}
  {\lccode`\~="2203 \lowercase{\global\catcode`~=\active \global\def~{\ensuremath{\exists}}}}
  {\lccode`\~="2212 \lowercase{\global\catcode`~=\active \global\def~{\ensuremath{-}}}}
  {\lccode`\~="2227 \lowercase{\global\catcode`~=\active \global\def~{\ensuremath{\wedge}}}}
  {\lccode`\~="2261 \lowercase{\global\catcode`~=\active \global\def~{\ensuremath{\equiv}}}}
  {\lccode`\~="2264 \lowercase{\global\catcode`~=\active \global\def~{\ensuremath{\leq}}}}
  {\lccode`\~="2265 \lowercase{\global\catcode`~=\active \global\def~{\ensuremath{\geq}}}}
  {\lccode`\~="27E8 \lowercase{\global\catcode`~=\active \global\def~{\ensuremath{\langle}}}}
  {\lccode`\~="27E9 \lowercase{\global\catcode`~=\active \global\def~{\ensuremath{\rangle}}}}
  {\lccode`\~="27FA \lowercase{\global\catcode`~=\active \global\def~{\ensuremath{\iff}}}}
\fi
\ifdefined\XeTeXversion
  \usepackage{xeCJK}
  \setCJKmainfont{SimSun}
\fi

\newtheorem{theorem}{Theorem}
\theoremstyle{definition}
\newtheorem{conjecture}{Conjecture}

\title{Interleaving structure and mod-2 pattern laws for the pendant-ladder
matching-count recurrences\\ (Lean 4 machine-checked claim of record)}
\author{[author name placeholder]\thanks{Contact and ORCID: placeholders by SOP
08b, to be filled by the author at upload; the pipeline writes no personal
information.}}
\date{September 28, 2026}

\begin{document}
\maketitle

\begin{abstract}
Let $a(n)$ be the number of matchings (independent edge sets, the empty one
included) of the ladder graph $P_n\,\square\,K_2$ with one pendant leaf attached
at every 1-based odd column on a uniform rail, starting
$3,10,46,141,660,2015,\dots$. Three integer sequences are defined by literal
seeds and linear recurrences: one by eight seeds and the even-lag order-$8$
recurrence $x(n)=16x(n-2)-25x(n-4)+10x(n-6)-x(n-8)$, and two subsequences by
four seeds each under the shared order-$4$ recurrence with characteristic
polynomial $x^4-16x^3+25x^2-10x+1$. We state and machine-check in Lean 4
(v4.34.0, mathlib v4.34.0) four laws for these recurrence-defined sequences:
the interleaving identity between the full sequence and the two subsequences;
an explicit period-$12$ mod-$2$ pattern of the full sequence; oddness of the
odd-position subsequence exactly at $k\equiv 0,4\pmod 6$; and evenness of the
even-position subsequence exactly at $k\equiv 0\pmod 3$. Every proof compiles
with no \texttt{sorry}, and \texttt{\#print axioms} reports only \{propext,
Quot.sound\}. The odd-position values coincide with OEIS A386889 (Dresden and
Demirkol, 2025-09-04), an occupied record; the novelty asserted here is limited
to the full order-$8$ recurrence, the even-position subsequence, the
interleaving identity and the mod-$2$ patterns. The combinatorial
identification of the sequences with matching counts is not kernel-proved and
is stated as a conjecture. This note is a priority claim of record, not a full
paper.
\end{abstract}

\section{Introduction}

Four theorems about three integer sequences defined by initial values and linear
recurrences are reported here; conjecturally the sequences count matchings of a
pendant-leaf ladder graph \cite{godsil-gutman}, and the recurrences were
obtained by exact rational Berlekamp--Massey fitting of enumerated data, then
independently recomputed \cite{massey}. The proofs were produced by the AI4Math
pipeline, in which LLM workers generate statements and proof scripts and the
Lean 4 kernel is the sole adjudicator at every gate: statements are frozen by
hash before proving, and a separate audit re-computes the symmetric difference,
scans the whole file for \texttt{sorry}/\texttt{admit} including comments, and
re-prints the axiom dependencies \cite{lean4,mathlib}. Verification strength is
stated below verbatim. This note is a priority claim of record deposited so that
a dated, citable public record exists ahead of any narrative paper; a fuller
paper for the same result is prepared separately and will reference this
deposit. No literature review is attempted here.

\section{Statement}\label{sec:statement}

The three sequences are \emph{defined} by the displayed initial values and
recurrences; $0$-based indexing with $\mathtt{apend}\,n=a(n+1)$,
$\mathtt{opend}\,k=a(2k+1)$ (odd-position side) and $\mathtt{epend}\,k=a(2k+2)$
(even-position side). Occupancy fence, carried verbatim from the audit record:
the numerical values of $\mathtt{opend}$ coincide with OEIS
\textbf{A386889} \cite{a386889} (Dresden and Demirkol, 2025-09-04; occupied,
values cross-checked to $k=15$); the novelty of this line is strictly limited
to the full order-$8$ recurrence plus the even-position subsequence
($\mathtt{epend}$ side); no ``new sequence'' claim about the whole family is
made anywhere in this note.

% LEAN: tasks/20260927-mossad-pend :: apend_interleave grade=完全证明
\begin{theorem}[interleaving identity]\label{thm:t1}
For the sequences of Listing~\ref{lst:s1} and every $k\in\mathbb{N}$,
$\mathtt{apend}(2k)=\mathtt{opend}(k)$ and
$\mathtt{apend}(2k+1)=\mathtt{epend}(k)$: the full order-$8$ recurrence is
exactly the interleaving of two order-$4$ subsequences sharing the
characteristic polynomial $x^4-16x^3+25x^2-10x+1$.
\end{theorem}

\begin{lstlisting}[caption={Frozen formal statements, verbatim lines 27--68 of
\texttt{formalized/01-pend-statements.lean} (sha256 \texttt{78bccc97}): the
three sequence definitions and Theorem~\ref{thm:t1}.},label={lst:s1}]
/-- **序列 apend（主列，order-8 仅偶滞后递推）**。
初值 a(0..7) = 3, 10, 46, 141, 660, 2015, 9440, 28814；
递推 a(n+8) = 16·a(n+6) − 25·a(n+4) + 10·a(n+2) − a(n)（n ≥ 0）。 -/
def apend : ℕ → ℤ
  | 0 => 3
  | 1 => 10
  | 2 => 46
  | 3 => 141
  | 4 => 660
  | 5 => 2015
  | 6 => 9440
  | 7 => 28814
  | n + 8 =>
      16 * apend (n + 6) - 25 * apend (n + 4) + 10 * apend (n + 2) - apend n

/-- **序列 opend（奇位侧子列，四阶递推）**。
初值 o(0..3) = 3, 46, 660, 9440（= apend 在下标 0,2,4,6 处之值；
数值与 OEIS A386889 一致——该侧已占位，见文件头围栏）；
递推 o(k+4) = 16·o(k+3) − 25·o(k+2) + 10·o(k+1) − o(k)（k ≥ 0）。 -/
def opend : ℕ → ℤ
  | 0 => 3
  | 1 => 46
  | 2 => 660
  | 3 => 9440
  | k + 4 =>
      16 * opend (k + 3) - 25 * opend (k + 2) + 10 * opend (k + 1) - opend k

/-- **序列 epend（偶位侧子列，四阶递推）**。
初值 e(0..3) = 10, 141, 2015, 28814（= apend 在下标 1,3,5,7 处之值）；
递推 e(k+4) = 16·e(k+3) − 25·e(k+2) + 10·e(k+1) − e(k)（k ≥ 0）。 -/
def epend : ℕ → ℤ
  | 0 => 10
  | 1 => 141
  | 2 => 2015
  | 3 => 28814
  | k + 4 =>
      16 * epend (k + 3) - 25 * epend (k + 2) + 10 * epend (k + 1) - epend k

/-- **T1（主攻）**：主列 order-8 递推 = 两条同系数四阶子列的交织——
对一切 k，apend (2k) = opend k 且 apend (2k+1) = epend k。 -/
theorem apend_interleave (k : ℕ) :
    apend (2 * k) = opend k ∧ apend (2 * k + 1) = epend k := by
\end{lstlisting}

% LEAN: tasks/20260927-mossad-pend :: apend_mod2_period12 grade=完全证明
\begin{theorem}[period-12 mod-2 pattern of the full sequence]\label{thm:t2}
For every $n\in\mathbb{N}$, $\mathtt{apend}(n)\bmod 2$ depends only on
$n\bmod 12$: it is $1$ exactly for $n\bmod 12\in\{0,3,5,8,9,11\}$ and $0$
otherwise.
\end{theorem}

\begin{lstlisting}[caption={Frozen formal statement, verbatim lines 71--89 of
\texttt{formalized/01-pend-statements.lean}: the pattern \texttt{pat2} and
Theorem~\ref{thm:t2}.},label={lst:s2}]
/-- **余数样式 pat2**（12 留数类 ↦ 余数）：
0↦1, 1↦0, 2↦0, 3↦1, 4↦0, 5↦1, 6↦0, 7↦0, 8↦1, 9↦1, 10↦0, 11↦1。 -/
def pat2 : ℕ → ℤ
  | 0 => 1
  | 1 => 0
  | 2 => 0
  | 3 => 1
  | 4 => 0
  | 5 => 1
  | 6 => 0
  | 7 => 0
  | 8 => 1
  | 9 => 1
  | 10 => 0
  | _ => 1

/-- **T2（伴随）**：apend 的 mod 2 余数仅依赖 n mod 12，样式为 pat2。 -/
theorem apend_mod2_period12 (n r : ℕ) (hr : n % 12 = r) :
    apend n % 2 = pat2 r := by
\end{lstlisting}

% LEAN: tasks/20260927-mossad-pend :: opend_mod2_period6 grade=完全证明
\begin{theorem}[oddness pattern of the odd-position subsequence]\label{thm:t3o}
For every $k\in\mathbb{N}$, $\mathtt{opend}(k)$ is odd if and only if
$k\equiv 0$ or $4\pmod 6$.
\end{theorem}

\begin{lstlisting}[caption={Frozen formal statements, verbatim lines 92--106 of
\texttt{formalized/01-pend-statements.lean}: the patterns \texttt{pat6} and
\texttt{pat3} and Theorem~\ref{thm:t3o}.},label={lst:s3}]
/-- **余数样式 pat6**：0↦1, 4↦1，其余↦0（订正后口径：奇 ⟺ k ≡ 0 或 4 (mod 6)）。 -/
def pat6 : ℕ → ℤ
  | 0 => 1
  | 4 => 1
  | _ => 0

/-- **余数样式 pat3**：0↦0，其余↦1（订正后口径：偶 ⟺ k ≡ 0 (mod 3)）。 -/
def pat3 : ℕ → ℤ
  | 0 => 0
  | _ => 1

/-- **T3 之 opend（伴随）**：opend k 的 mod 2 余数仅依赖 k mod 6，样式为 pat6
（即 opend k 为奇 ⟺ k ≡ 0 或 4 (mod 6)）。 -/
theorem opend_mod2_period6 (k r : ℕ) (hr : k % 6 = r) :
    opend k % 2 = pat6 r := by
\end{lstlisting}

% LEAN: tasks/20260927-mossad-pend :: epend_mod2_period3 grade=完全证明
\begin{theorem}[evenness pattern of the even-position subsequence]\label{thm:t3e}
For every $k\in\mathbb{N}$, $\mathtt{epend}(k)$ is even if and only if
$k\equiv 0\pmod 3$.
\end{theorem}

\begin{lstlisting}[caption={Frozen formal statement, verbatim lines 109--112 of
\texttt{formalized/01-pend-statements.lean}: Theorem~\ref{thm:t3e}.},
label={lst:s4}]
/-- **T3 之 epend（伴随）**：epend k 的 mod 2 余数仅依赖 k mod 3，样式为 pat3
（即 epend k 为偶 ⟺ k ≡ 0 (mod 3)）。 -/
theorem epend_mod2_period3 (k r : ℕ) (hr : k % 3 = r) :
    epend k % 2 = pat3 r := by
\end{lstlisting}

\begin{conjecture}[combinatorial bridge, not proved here]\label{conj:bridge}
For all $n\ge 0$, $\mathtt{apend}(n)$ equals the number of matchings of the
ladder graph $P_{n+1}\,\square\,K_2$ with one pendant leaf attached at every
1-based odd column on a uniform rail (the alternating-attachment variant
differs from the uniform-side convention from $n=3$ on and is a different
count).
\end{conjecture}

\section{Proof}\label{sec:proof}

All four proofs are strong induction on the index with the seed window
discharged by computation and the step closed by recurrence unfolding and
residue tables (\texttt{interval\_cases}/\texttt{decide},
\texttt{mod\_cases}, \texttt{omega}); Theorems~\ref{thm:t2}, \ref{thm:t3o} and
\ref{thm:t3e} first prove a clean main lemma without the residue hypothesis
and then substitute it back. The complete final file is reproduced verbatim
below ($259$ lines, no excerpting); it also ships in this deposit under
\texttt{proofs/}.

\begin{lstlisting}[basicstyle=\ttfamily\footnotesize,caption={Complete proofs of
Theorems~\ref{thm:t1}--\ref{thm:t3e}, verbatim lines 1--259 (whole file) of
\texttt{final/01-pend-proved.lean} (sha256 \texttt{be3f31cb}); the smaller
monospace width is typesetting only, listing bytes are the frozen file's own.},
label={lst:proof}]
/-
  AI4Math 流水线 · AI 生成 · 2026-09-27（终稿生成 2026-09-28）
  部门：02 形式化部（dept-formalize）statement 冻结 + 04 军政部（dept-prove）proof body 终稿
  任务线：mossad-pend（pool-mossad-01，pool Part A；三线兄弟线之 pend 线）
  来源卡：tasks/20260927-mossad-pend/card.md（闸门一 2026-09-27 立项）
  主张定稿：tasks/20260927-mossad-pend/phase0/claims.md（T1 主攻 + T2/T3 伴随；T4 备选不进本批）
  初值数据出处：phase0/counts.txt（双算法互证 PASS）+ claims.md 主代理复核订正节
    （order-8 递推、两四阶子列递推、T1 交织恒等式、T2 样式在 n≤30 数据窗口全项 PASS）；
    本文件五条 def 的初值/系数/样式逐字照抄 claims.md Lean 草稿（T3 取订正后形态，
    格式上仅把每行多分枝改为每行一分枝，语义不变）。

  口径（按任务卡锁定）：本文件（含注释）只使用「初值 + 整数系数递推定义的纯序列」
  语言（ℕ → ℤ），命题只谈序列的交织/模余数性质；组合对象语义一律不进本层。
  索引口径（0-based Lean ↔ 1-based 计数）：
  apend n = a(n+1)；opend k = a(2k+1)（奇位侧）；epend k = a(2k+2)（偶位侧）。

  【占位围栏】opend 的数值 = OEIS A386889（Dresden & Demirkol，2025-09-04 挂名，
  已占位，数值核到 k=15）；本线新颖性主张严格限定「完整 order-8 递推 +
  偶数子列（epend 侧）」；禁止任何「新序列」泛写。statement 层仅为纯序列命题，
  占位状态不改变命题真伪，仅约束下游报告/论文措辞。

  proof body 由军政部完成；良定义门由宪兵 gate-batch 集中跑。
  （来源探针：T1=attempts/r6-t1.lean、T2=attempts/r5-t2.lean、
  T3o=attempts/r4-t3o.lean、T3e=attempts/r3-step-t3e.lean，均 EXIT=0 通过。）
-/

import Mathlib

/-- **序列 apend（主列，order-8 仅偶滞后递推）**。
初值 a(0..7) = 3, 10, 46, 141, 660, 2015, 9440, 28814；
递推 a(n+8) = 16·a(n+6) − 25·a(n+4) + 10·a(n+2) − a(n)（n ≥ 0）。 -/
def apend : ℕ → ℤ
  | 0 => 3
  | 1 => 10
  | 2 => 46
  | 3 => 141
  | 4 => 660
  | 5 => 2015
  | 6 => 9440
  | 7 => 28814
  | n + 8 =>
      16 * apend (n + 6) - 25 * apend (n + 4) + 10 * apend (n + 2) - apend n

/-- **序列 opend（奇位侧子列，四阶递推）**。
初值 o(0..3) = 3, 46, 660, 9440（= apend 在下标 0,2,4,6 处之值；
数值与 OEIS A386889 一致——该侧已占位，见文件头围栏）；
递推 o(k+4) = 16·o(k+3) − 25·o(k+2) + 10·o(k+1) − o(k)（k ≥ 0）。 -/
def opend : ℕ → ℤ
  | 0 => 3
  | 1 => 46
  | 2 => 660
  | 3 => 9440
  | k + 4 =>
      16 * opend (k + 3) - 25 * opend (k + 2) + 10 * opend (k + 1) - opend k

/-- **序列 epend（偶位侧子列，四阶递推）**。
初值 e(0..3) = 10, 141, 2015, 28814（= apend 在下标 1,3,5,7 处之值）；
递推 e(k+4) = 16·e(k+3) − 25·e(k+2) + 10·e(k+1) − e(k)（k ≥ 0）。 -/
def epend : ℕ → ℤ
  | 0 => 10
  | 1 => 141
  | 2 => 2015
  | 3 => 28814
  | k + 4 =>
      16 * epend (k + 3) - 25 * epend (k + 2) + 10 * epend (k + 1) - epend k

/-- **T1（主攻）**：主列 order-8 递推 = 两条同系数四阶子列的交织——
对一切 k，apend (2k) = opend k 且 apend (2k+1) = epend k。 -/
theorem apend_interleave (k : ℕ) :
    apend (2 * k) = opend k ∧ apend (2 * k + 1) = epend k := by
  induction k using Nat.strong_induction_on with
  | h k ih =>
    by_cases hk : k < 4
    · interval_cases k <;> decide
    · obtain ⟨j, rfl⟩ : ∃ j, k = j + 4 := ⟨k - 4, by omega⟩
      constructor
      · have e8 : 2 * (j + 4) = (2 * j) + 8 := by ring
        rw [e8]
        simp only [apend, opend]
        rw [show 2 * j + 6 = 2 * (j + 3) by ring,
            show 2 * j + 4 = 2 * (j + 2) by ring,
            show 2 * j + 2 = 2 * (j + 1) by ring,
            (ih (j + 3) (by omega)).1, (ih (j + 2) (by omega)).1,
            (ih (j + 1) (by omega)).1, (ih j (by omega)).1]
      · have e9 : 2 * (j + 4) + 1 = (2 * j + 1) + 8 := by ring
        rw [e9]
        simp only [apend, epend]
        rw [show 2 * j + 1 + 6 = 2 * (j + 3) + 1 by ring,
            show 2 * j + 1 + 4 = 2 * (j + 2) + 1 by ring,
            show 2 * j + 1 + 2 = 2 * (j + 1) + 1 by ring,
            (ih (j + 3) (by omega)).2, (ih (j + 2) (by omega)).2,
            (ih (j + 1) (by omega)).2, (ih j (by omega)).2]

/-- **余数样式 pat2**（12 留数类 ↦ 余数）：
0↦1, 1↦0, 2↦0, 3↦1, 4↦0, 5↦1, 6↦0, 7↦0, 8↦1, 9↦1, 10↦0, 11↦1。 -/
def pat2 : ℕ → ℤ
  | 0 => 1
  | 1 => 0
  | 2 => 0
  | 3 => 1
  | 4 => 0
  | 5 => 1
  | 6 => 0
  | 7 => 0
  | 8 => 1
  | 9 => 1
  | 10 => 0
  | _ => 1

/-- **T2（伴随）**：apend 的 mod 2 余数仅依赖 n mod 12，样式为 pat2。 -/
theorem apend_mod2_period12 (n r : ℕ) (hr : n % 12 = r) :
    apend n % 2 = pat2 r := by
  have main : ∀ n : ℕ, apend n % 2 = pat2 (n % 12) := by
    intro n
    induction n using Nat.strong_induction_on with
    | h n ih =>
      by_cases hn : n < 8
      · interval_cases n <;> decide
      · obtain ⟨j, rfl⟩ : ∃ j, n = j + 8 := ⟨n - 8, by omega⟩
        have h1 : apend (j + 8) % 2 = (apend (j + 4) % 2 + apend j % 2) % 2 := by
          simp only [apend]
          omega
        rw [h1, ih (j + 4) (by omega), ih j (by omega)]
        mod_cases hj : j % 12
        · have hj0 : j % 12 = 0 := by simpa [Nat.ModEq] using hj
          have h4 : (j + 4) % 12 = 4 := by omega
          have h8 : (j + 8) % 12 = 8 := by omega
          rw [hj0, h4, h8]; decide
        · have hj1 : j % 12 = 1 := by simpa [Nat.ModEq] using hj
          have h4 : (j + 4) % 12 = 5 := by omega
          have h8 : (j + 8) % 12 = 9 := by omega
          rw [hj1, h4, h8]; decide
        · have hj2 : j % 12 = 2 := by simpa [Nat.ModEq] using hj
          have h4 : (j + 4) % 12 = 6 := by omega
          have h8 : (j + 8) % 12 = 10 := by omega
          rw [hj2, h4, h8]; decide
        · have hj3 : j % 12 = 3 := by simpa [Nat.ModEq] using hj
          have h4 : (j + 4) % 12 = 7 := by omega
          have h8 : (j + 8) % 12 = 11 := by omega
          rw [hj3, h4, h8]; decide
        · have hj4 : j % 12 = 4 := by simpa [Nat.ModEq] using hj
          have h4 : (j + 4) % 12 = 8 := by omega
          have h8 : (j + 8) % 12 = 0 := by omega
          rw [hj4, h4, h8]; decide
        · have hj5 : j % 12 = 5 := by simpa [Nat.ModEq] using hj
          have h4 : (j + 4) % 12 = 9 := by omega
          have h8 : (j + 8) % 12 = 1 := by omega
          rw [hj5, h4, h8]; decide
        · have hj6 : j % 12 = 6 := by simpa [Nat.ModEq] using hj
          have h4 : (j + 4) % 12 = 10 := by omega
          have h8 : (j + 8) % 12 = 2 := by omega
          rw [hj6, h4, h8]; decide
        · have hj7 : j % 12 = 7 := by simpa [Nat.ModEq] using hj
          have h4 : (j + 4) % 12 = 11 := by omega
          have h8 : (j + 8) % 12 = 3 := by omega
          rw [hj7, h4, h8]; decide
        · have hj8 : j % 12 = 8 := by simpa [Nat.ModEq] using hj
          have h4 : (j + 4) % 12 = 0 := by omega
          have h8 : (j + 8) % 12 = 4 := by omega
          rw [hj8, h4, h8]; decide
        · have hj9 : j % 12 = 9 := by simpa [Nat.ModEq] using hj
          have h4 : (j + 4) % 12 = 1 := by omega
          have h8 : (j + 8) % 12 = 5 := by omega
          rw [hj9, h4, h8]; decide
        · have hj10 : j % 12 = 10 := by simpa [Nat.ModEq] using hj
          have h4 : (j + 4) % 12 = 2 := by omega
          have h8 : (j + 8) % 12 = 6 := by omega
          rw [hj10, h4, h8]; decide
        · have hj11 : j % 12 = 11 := by simpa [Nat.ModEq] using hj
          have h4 : (j + 4) % 12 = 3 := by omega
          have h8 : (j + 8) % 12 = 7 := by omega
          rw [hj11, h4, h8]; decide
  rw [← hr]; exact main n

/-- **余数样式 pat6**：0↦1, 4↦1，其余↦0（订正后口径：奇 ⟺ k ≡ 0 或 4 (mod 6)）。 -/
def pat6 : ℕ → ℤ
  | 0 => 1
  | 4 => 1
  | _ => 0

/-- **余数样式 pat3**：0↦0，其余↦1（订正后口径：偶 ⟺ k ≡ 0 (mod 3)）。 -/
def pat3 : ℕ → ℤ
  | 0 => 0
  | _ => 1

/-- **T3 之 opend（伴随）**：opend k 的 mod 2 余数仅依赖 k mod 6，样式为 pat6
（即 opend k 为奇 ⟺ k ≡ 0 或 4 (mod 6)）。 -/
theorem opend_mod2_period6 (k r : ℕ) (hr : k % 6 = r) :
    opend k % 2 = pat6 r := by
  have main : ∀ k : ℕ, opend k % 2 = pat6 (k % 6) := by
    intro k
    induction k using Nat.strong_induction_on with
    | h k ih =>
      by_cases hk : k < 4
      · interval_cases k <;> decide
      · obtain ⟨j, rfl⟩ : ∃ j, k = j + 4 := ⟨k - 4, by omega⟩
        have h1 : opend (j + 4) % 2 = (opend (j + 2) % 2 + opend j % 2) % 2 := by
          simp only [opend]
          omega
        rw [h1, ih (j + 2) (by omega), ih j (by omega)]
        mod_cases hj : j % 6
        · have hj0 : j % 6 = 0 := by simpa [Nat.ModEq] using hj
          have h2 : (j + 2) % 6 = 2 := by omega
          have h4 : (j + 4) % 6 = 4 := by omega
          rw [hj0, h2, h4]; decide
        · have hj1 : j % 6 = 1 := by simpa [Nat.ModEq] using hj
          have h2 : (j + 2) % 6 = 3 := by omega
          have h4 : (j + 4) % 6 = 5 := by omega
          rw [hj1, h2, h4]; decide
        · have hj2 : j % 6 = 2 := by simpa [Nat.ModEq] using hj
          have h2 : (j + 2) % 6 = 4 := by omega
          have h4 : (j + 4) % 6 = 0 := by omega
          rw [hj2, h2, h4]; decide
        · have hj3 : j % 6 = 3 := by simpa [Nat.ModEq] using hj
          have h2 : (j + 2) % 6 = 5 := by omega
          have h4 : (j + 4) % 6 = 1 := by omega
          rw [hj3, h2, h4]; decide
        · have hj4 : j % 6 = 4 := by simpa [Nat.ModEq] using hj
          have h2 : (j + 2) % 6 = 0 := by omega
          have h4 : (j + 4) % 6 = 2 := by omega
          rw [hj4, h2, h4]; decide
        · have hj5 : j % 6 = 5 := by simpa [Nat.ModEq] using hj
          have h2 : (j + 2) % 6 = 1 := by omega
          have h4 : (j + 4) % 6 = 3 := by omega
          rw [hj5, h2, h4]; decide
  rw [← hr]; exact main k

/-- **T3 之 epend（伴随）**：epend k 的 mod 2 余数仅依赖 k mod 3，样式为 pat3
（即 epend k 为偶 ⟺ k ≡ 0 (mod 3)）。 -/
theorem epend_mod2_period3 (k r : ℕ) (hr : k % 3 = r) :
    epend k % 2 = pat3 r := by
  have main : ∀ k : ℕ, epend k % 2 = pat3 (k % 3) := by
    intro k
    induction k using Nat.strong_induction_on with
    | h k ih =>
      by_cases hk : k < 4
      · interval_cases k <;> decide
      · obtain ⟨j, rfl⟩ : ∃ j, k = j + 4 := ⟨k - 4, by omega⟩
        have h1 : epend (j + 4) % 2 = (epend (j + 2) % 2 + epend j % 2) % 2 := by
          simp only [epend]
          omega
        rw [h1, ih (j + 2) (by omega), ih j (by omega)]
        mod_cases hj : j % 3
        · have hj0 : j % 3 = 0 := by simpa [Nat.ModEq] using hj
          have h2 : (j + 2) % 3 = 2 := by omega
          have h4 : (j + 4) % 3 = 1 := by omega
          rw [hj0, h2, h4]
          decide
        · have hj1 : j % 3 = 1 := by simpa [Nat.ModEq] using hj
          have h2 : (j + 2) % 3 = 0 := by omega
          have h4 : (j + 4) % 3 = 2 := by omega
          rw [hj1, h2, h4]
          decide
        · have hj2 : j % 3 = 2 := by simpa [Nat.ModEq] using hj
          have h2 : (j + 2) % 3 = 1 := by omega
          have h4 : (j + 4) % 3 = 0 := by omega
          rw [hj2, h2, h4]
          decide
  rw [← hr]; exact main k
\end{lstlisting}

\section{Verification evidence}

\begin{itemize}
\item Toolchain: Lean \texttt{v4.34.0} (\texttt{leanprover/lean4:v4.34.0}),
mathlib4 \texttt{v4.34.0}, revision \texttt{5ed29652}. Reproduction: with that
pin, \texttt{lake env lean} on \texttt{proofs/01-pend-proved.lean} compiles the
file with exit code $0$ and empty output (measured $70.0$\,s in the audit's
independent strict re-run).
\item Statement integrity: the frozen statement file\newline
\texttt{formalized/01-pend-statements.lean} ($113$ lines) has sha256
\begin{center}
\texttt{78bccc97afba79adea9658a2158d17b2}\allowbreak
\texttt{ed247b584d9f28266e6fc7914ba9f922},
\end{center}
and the final file \texttt{final/01-pend-proved.lean} ($259$ lines) has sha256
\begin{center}
\texttt{be3f31cb365aa7b805dce0d29e562b5d}\allowbreak
\texttt{fdb12617007f888e67564dacfd22540b}.
\end{center}
The audit extracted the ten declaration blocks (six definitions, four theorems,
docstrings included) from both files programmatically: $10/10$ byte-identical;
a whole-file diff shows exactly six hunks, namely three header-comment lines and
the four \texttt{sorry} placeholder lines replaced by proof bodies -- the
statement body itself has zero drift.
\item Kernel check: compiles with $0$ \texttt{sorry} (whole-file scan including
comments, case-insensitive, also covering \texttt{admit}, \texttt{sorryAx},
\texttt{native\_decide}, \texttt{axiom} and \texttt{unsafe}, all zero hits);
\texttt{\#print axioms} output, verbatim (kept literal; the smaller monospace
width is typesetting only):
\begin{lstlisting}[basicstyle=\ttfamily\footnotesize,frame=none]
'apend_interleave' depends on axioms: [propext, Quot.sound]
'apend_mod2_period12' depends on axioms: [propext, Quot.sound]
'opend_mod2_period6' depends on axioms: [propext, Quot.sound]
'epend_mod2_period3' depends on axioms: [propext, Quot.sound]
\end{lstlisting}
Each line is a subset of the whitelist \{propext, Quot.sound,
Classical.choice\} -- in fact Classical.choice is not even used. The strict
re-compile emits no error and no warning at all (stdout empty); the audit's
independent re-runs ($70.0$\,s strict, $41.3$\,s axioms) are byte-identical to
the self-reported outputs. The evidence ships as
\texttt{audit/audit-strict-compile.log} and
\texttt{audit/audit-axioms-sidecar.log} (the four lines above are that log's
payload; the shipped copy adds a three-line provenance header, and the audit's
re-run \texttt{audit/axioms-recheck.out} is byte-identical to the task's own
axiom printout).
\item Grading by the audit department: all four theorems \emph{complete proof};
lowest grade reached for the case = \emph{complete proof}; no scaffolding route
was used (the proof skeleton replicates a same-pool kernel-verified precedent,
which the audit ruled is not a scaffolding downgrade).
\item Audit chain (originals in this deposit under \texttt{audit/}): final
adjudication \texttt{final-01-pend.txt} (eight independent checks),
well-definedness verdict \texttt{welldef-verdict.md} with gate logs
\texttt{gate-t1.txt}, \texttt{gate-t2.txt}, \texttt{gate-t3o.txt},
\texttt{gate-t3e.txt} and the supplementary \texttt{ring} probe
\texttt{gate-t1-ring.txt} (degeneracy probes $33/33$ FAIL, i.e.\ no one-tactic
shortcut), the statement symmetric-difference record
\texttt{statement-diff.out}, the audit's re-verification
\texttt{axioms-recheck.out}, and this note's claim check
\texttt{claim-check.md} --- added to \texttt{audit/} after the audit department
returns its adjudication, and a prerequisite for publication. The
semantic-fidelity check (roundtrip re-translation) of the underlying task ran
on the same model node as its formalisation; that independence weakness is
mitigated, not removed, by kernel adjudication and by two independent audit
re-runs.
\end{itemize}

\section{Novelty evidence}\label{sec:novelty}

Adjudication copied from the pipeline's exclusivity review (C9 gate), which ran
before the problem was accepted and re-derived the data independently: verdict
\textbf{no occupying record found for the main sequence; the odd-position
subsequence is already occupied (OEIS A386889)}, value tier \textbf{new
sequence / new recurrence (with an occupancy warning)}. It is not a claim of
new mathematics and not a first-discovery claim, and no ``new sequence'' claim
about the whole family is made. Summary by channel; the full query log,
including every zero-hit query, is in this deposit at
\texttt{audit/c9-record.md}.

\begin{center}
\resizebox{\textwidth}{!}{%
\begin{tabular}{lll}
\hline
channel & queries & result \\
\hline
OEIS main-sequence variants \cite{oeis} & 11 & zero hits: original string,
drop-1/2/3 (4), $-1$, $+1$, $-1$ then drop-1, $+1$ then drop-1 (4), $-2$ (1),
$\times2$, $\times2$ then drop-1 (2); $\div2$ not applicable (first term odd),
recorded \\
OEIS odd-position subsequence & 6 & \textbf{A386889} \cite{a386889} on the
plain string and two reformulations; $-1$, $+1$, $\times2$ variants zero hits \\
OEIS even-position subsequence & 6 & zero hits (plain, drop-1, 8-term window,
$-1$, $+1$, $\times2$) \\
OEIS keywords & 4 & unrelated noise only (A000045, A011973) \\
zbMATH & 6 & two clean zero readings for pendant-ladder matching queries;
Farrell 1978 \emph{Matchings in ladders} and others concern plain ladders, no
pendant structures \\
OpenAlex & 3 & 235 full-text-noise hits; 0 hits; one query HTTP 429, recorded
as unverified \\
arXiv & 5 & three exact-phrase zeros; author searches found no tiling paper
behind A386889 \\
Math Stack Exchange & 1 & zero items \\
general web layer & 13 attempts & all unverified (channel failure, recorded as
such, not used as zero-literature evidence) \\
library layer & 3 & mathlib has matching predicates but no matching-count
theorem; compfiles none; sequencelib none \\
public Lean ecosystem & 4 & repository searches and one full-tree scan: zero
hits \\
\hline
\end{tabular}%
}
\end{center}

Three points of the record are part of the evidence. First, the review's anchor
self-proof (Fibonacci and a known tatami sequence returned their correct OEIS
entries) fixes what ``zero hit'' means on the sequence channel. Second, the
occupying entry: OEIS A386889 (Dresden and Demirkol, 2025-09-04) counts
$1\times1$/$1\times2$ tilings of a $3\times(2n-1)$ strip with every other
top-row cell removed; the review aligned 13 terms of the odd-position
subsequence against it offset-for-offset and independently verified the graph
isomorphism with the pendant ladder (degree sequences and an explicit bijection
at $n=2$), and its recurrence signature $(16,-25,10,-1)$ equals the one fitted
here --- so the odd-position side is \emph{occupied} and carries the fence
stated in Section~\ref{sec:statement}. Third, the externally supplied review
had queried only six main-sequence variants and missed exactly this derived
occupancy; the degradation notes (OpenAlex 429, general web layer unavailable,
no Chinese-language database searched) are on file. If any channel later
returns a hit on the remaining novelty (full order-$8$ recurrence,
even-position subsequence, interleaving, mod-2 patterns), the tier above
weakens accordingly, and this deposit will be amended by a later version rather
than by editing this record.

\section{Scope and limitations}

Grade and boundary, stated as the audit states them. (a) The theorems cover the
recurrence-defined sequences only; their equality with pendant-ladder matching
counts is not kernel-proved and is Conjecture~\ref{conj:bridge}. Its support is
OEIS/literature endorsement plus local numeric probes: exact counting from the
graph definition by three independent methods agreeing on the seed window and
extended to $n\le 30$; the odd-position side was aligned with A386889 over 13
terms and cross-checked to $k=15$; the well-definedness audit independently
recomputed all statements on the window $n\le 30$ / $k\le 14$. The enumerators
themselves are not verified, and the scripts and their outputs stay in the
working repository under \texttt{tasks/20260927-mossad-pend/}. (b)
Theorems~\ref{thm:t3o} and \ref{thm:t3e} are stated in their corrected forms
(odd iff $k\equiv 0,4 \pmod 6$; even iff $k\equiv 0 \pmod 3$); the earlier
off-by-one formulations ($k\equiv 1,5$; $k\equiv 1$) are deprecated and must
not be cited. (c) Indexing is $0$-based: $\mathtt{apend}\,n=a(n+1)$,
$\mathtt{opend}\,k=a(2k+1)$, $\mathtt{epend}\,k=a(2k+2)$. (d) The recurrences
are computational findings: Berlekamp--Massey fits over $\mathbb{Q}$ verified
digit-for-digit against enumeration and independently recomputed, not derived
from a transfer-matrix argument. The older ``alternating attachment is robust''
belief is refuted: alternating and uniform-side attachments differ from $n=3$
on ($3,10,47,143,\dots$ vs $3,10,46,141,\dots$, three methods agreeing), so the
combinatorial reading is fixed to the uniform-side convention. (e) The word
\emph{minimal} in ``minimal period $12/6/3$'' is supported by the data window,
the audit's independent recomputation ruling out the candidate divisors, and a
GF(2) factorisation cross-check; the kernel proves the universal statements
themselves but not minimality. (f) The gate-four sign-off of the underlying
task was exercised by the author's instruction to deposit this claim on
2026-09-28, while the task report file still carries its draft header; both
facts are recorded here as they stand. (g) A mod-4 pattern and
arbitrary-modulus period questions were left as an unattempted option and are
not implied anywhere. (h) This note carries no literature review and is a
deposit record, not a survey; the narrative paper for the same result is
separate.

\section*{Data availability}

The deposit package for this claim contains the claim note (LaTeX source and
PDF), the Lean sources (\texttt{proofs/}: frozen statement file, final proof
file, \texttt{lean-toolchain}), the verification and audit logs, and the full
C9 query log with zero-hit entries (\texttt{audit/}). Zenodo: the DOI is
reserved at upload by the author and is recorded in \texttt{metadata/README.md}
of this package once created; see \texttt{UPLOAD.md} in the source repository.
Claim note under CC BY 4.0, code under MIT.

\section*{Use of generative AI}

(a) AI performed: candidate generation and the exclusivity review,
autoformalisation of the six definitions and four statements, the roundtrip
semantic check, proof search and proof-script writing, error classification and
repair, and the assembly of this note including its verbatim listings. (b)
Model, node and budget: all task-stage sampling ran on one node, wire-id
\texttt{anthropic/a6api-main/kimi-k3}, with no cross-node switching, and
consumed $2$ LLM calls of a permitted $8$ (roundtrip adjudication: one call
truncated at the token cap and retried on the same node, one succeeded) and
$20$ kernel compiles of a permitted $24$, copied verbatim from the task ledger.
This note consumed $0$ pipeline LLM calls against a cap of $4$; drafting was
mechanical assembly. (c) Final adjudication: all statements and proofs reported
here were checked by the Lean 4 kernel: every proof compiles with no
\texttt{sorry}, and \texttt{\#print axioms} reports only \{propext,
Quot.sound\}. (d) The AI system is not an author; the human author reviewed the
evidence and is responsible for all content.

\section*{Author contributions}

CRediT-style: the human author adjudicated topic selection, statement semantics
and the deposit decision; the AI4Math pipeline (an LLM) generated the candidate
propositions, the proof searches and the text drafts. The AI system is not
listed as an author.

\begin{thebibliography}{9}
% Each entry verified on 2026-09-28 against Crossref title/page metadata (the
% four DOIs) and live fetches of oeis.org and the A386889 page. Raw responses:
% claims/pendant-ladder-interleave/audit/bib-verify/. No entry is quoted from
% memory; the Lean 4 DOI suffix is _37 (the _27 suffix resolves to a Collatz
% paper in the same volume -- a template defect documented in earlier packages).
\bibitem{lean4}
L.~de Moura and S.~Ullrich.
\newblock The Lean 4 theorem prover and programming language.
\newblock \emph{Automated Deduction -- CADE-28}, LNCS, pp.~625--635, 2021.
\url{https://doi.org/10.1007/978-3-030-79876-5_37}

\bibitem{mathlib}
The mathlib Community.
\newblock The Lean mathematical library.
\newblock \emph{CPP 2020}, pp.~367--381, 2020.
\url{https://doi.org/10.1145/3372885.3373824}

\bibitem{massey}
J.~L. Massey.
\newblock Shift-register synthesis and BCH decoding.
\newblock \emph{IEEE Transactions on Information Theory} 15 (1969),
pp.~122--127.
\url{https://doi.org/10.1109/TIT.1969.1054260}

\bibitem{godsil-gutman}
C.~D. Godsil and I.~Gutman.
\newblock On the theory of the matching polynomial.
\newblock \emph{Journal of Graph Theory} 5 (1981), pp.~137--144.
\url{https://doi.org/10.1002/jgt.3190050203}

\bibitem{oeis}
The OEIS Foundation.
\newblock \emph{The On-Line Encyclopedia of Integer Sequences}.
\newblock \url{https://oeis.org/}

\bibitem{a386889}
G.~Dresden and D.~Demirkol.
\newblock Number of ways to use $1\times1$ squares and $1\times2$ dominos (in
any orientation) to tile a $3\times(2n-1)$ strip with every other cell in the
top row removed.
\newblock \emph{The On-Line Encyclopedia of Integer Sequences}, entry A386889,
2025-09-04.
\url{https://oeis.org/A386889}
\end{thebibliography}

\end{document}
